% =========================================================
% Algebraic Structures — Operations and Their Laws
% =========================================================
% Dependency chain: sets → functions → operations (this file).
% An n-ary operation is a function whose codomain is the same set
% as its domain factors. Closure is built into the definition.
% Order: general definition → nullary → unary → binary → laws → distinguished elements.

\section{Operations}

% ---------------------------------------------------------
% TOOLKIT — Operations
% ---------------------------------------------------------
\begin{tcolorbox}[colback=gray!6, colframe=gray!40, arc=2pt,
  left=6pt, right=6pt, top=4pt, bottom=4pt,
  title={\small\textbf{Operations — Quick Reference}},
  fonttitle=\small\bfseries]
\small
\begin{tabular}{l l l}
\toprule
\textbf{Concept} & \textbf{Definition} & \textbf{Detail} \\
\midrule
$n$-ary operation  & $\omega : A^n \to A$                           & \hyperref[def:nary-op]{↓} \\
Nullary operation  & $\omega : A^0 \to A$; selects an element of $A$         & \hyperref[def:nullary-op]{↓} \\
Unary operation    & $f : A \to A$                                  & \hyperref[def:unary-op]{↓} \\
Binary operation   & $\star : A \times A \to A$                     & \hyperref[def:binary-op]{↓} \\
Closure            & Consequence of $\omega : A^n \to A$; outputs stay in $A$ & \hyperref[def:binary-op]{↓} \\
\bottomrule
\end{tabular}
\end{tcolorbox}

\vspace{1em}

% ---------------------------------------------------------
% n-ary operation
% ---------------------------------------------------------

\begin{tcolorbox}[colback=defbox, colframe=defborder, arc=2pt,
  left=6pt, right=6pt, top=4pt, bottom=4pt,
  title={\small\textbf{Definition ($n$-ary Operation)}},
  fonttitle=\small\bfseries]
\label{def:nary-op}
Let $A$ be a set and $n \in \mathbb{N}$.
An \emph{$n$-ary operation} on $A$ is a function
\[
  \omega : A^n \to A.
\]
We call $A$ the \emph{underlying set} and $n$ the \emph{arity} of $\omega$.
\end{tcolorbox}

\begin{remark}[Fully quantified form]
$\omega$ is an $n$-ary operation on $A$
$\;\iff\;$
$(\forall a_1, \ldots, a_n \in A)\;
(\exists!\, b \in A)\;
\omega(a_1, \ldots, a_n) = b$.
\end{remark}

\begin{remark}[Operations are special functions]
An $n$-ary operation is a function $\omega : A^n \to A$ with one
distinguishing feature: the codomain is the \emph{same set} as the
domain factors.
This is what separates operations from arbitrary functions.
A function $f : A \to B$ with $B \neq A$ moves elements out of $A$;
an operation $\omega : A^n \to A$ stays within $A$.
\end{remark}

\begin{remark}[Closure is definitional]
The condition that $\omega(a_1, \ldots, a_n) \in A$ for all
$a_i \in A$ is not a property to be checked — it is built into the
type signature $\omega : A^n \to A$.
In older presentations, closure is listed as an axiom for groups and rings.
Here it is absorbed into the definition of operation itself.
\end{remark}

\begin{remark}[Examples beyond arithmetic]
Operations need not involve numbers.
On the set of plane symmetries of a square, composition of symmetries is a
binary operation.
On the set of strings over an alphabet, concatenation is a binary operation.
On the power set $\mathcal{P}(A)$, union and intersection are binary
operations.
On a vector space, scalar multiplication
$\mathbb{F} \times V \to V$ is a binary operation with a mixed domain
— a generalisation of the above that some authors call an \emph{external} operation.
\end{remark}

% ---------------------------------------------------------
% Nullary operation
% ---------------------------------------------------------

\begin{tcolorbox}[colback=defbox, colframe=defborder, arc=2pt,
  left=6pt, right=6pt, top=4pt, bottom=4pt,
  title={\small\textbf{Definition (Nullary Operation / Distinguished Element)}},
  fonttitle=\small\bfseries]
\label{def:nullary-op}
A \emph{nullary operation} on $A$ is a function $\omega : A^0 \to A$,
where $A^0 = \{\emptyset\}$ is the one-element set (the empty Cartesian
product).
Such an operation selects exactly one element $e = \omega(\emptyset) \in A$,
called a \emph{distinguished element} or \emph{constant} of $A$.
\end{tcolorbox}

\begin{remark}[Fully quantified form]
A nullary operation on $A$ is exactly a choice of a single element
$e \in A$.
The function perspective: $\omega : \{\emptyset\} \to A$,
$\omega(\emptyset) = e$.
\end{remark}

\begin{remark}[Why nullary operations matter]
In model-theoretic signatures, identity elements, zero elements, and
other distinguished constants are nullary operations — they are part of the
signature, not derived objects.
Writing ``the group $(G, \cdot, e, {}^{-1})$'' names a binary operation
$\cdot$, a nullary operation $e$ (the identity), and a unary operation
${}^{-1}$ (the inverse map).
This is why the identity is built into the signature rather than
existentially quantified in an axiom: it makes the structure a function
symbol, not a property.
\end{remark}

% ---------------------------------------------------------
% Unary operation
% ---------------------------------------------------------

\begin{tcolorbox}[colback=defbox, colframe=defborder, arc=2pt,
  left=6pt, right=6pt, top=4pt, bottom=4pt,
  title={\small\textbf{Definition (Unary Operation)}},
  fonttitle=\small\bfseries]
\label{def:unary-op}
A \emph{unary operation} on $A$ is a function
\[
  f : A \to A.
\]
\end{tcolorbox}

\begin{remark}[Fully quantified form]
$f$ is a unary operation on $A$
$\;\iff\;$
$(\forall a \in A)\;(\exists!\, b \in A)\; f(a) = b$.
\end{remark}

\begin{remark}[Examples]
Negation $a \mapsto -a$ on $\mathbb{Z}$ is a unary operation.
Multiplicative inverse $a \mapsto a^{-1}$ on $\mathbb{R} \setminus \{0\}$
is a unary operation.
Complex conjugation $z \mapsto \bar{z}$ on $\mathbb{C}$ is a unary operation.
In each case the output is a new element of the \emph{same} set.
\end{remark}

\begin{tcolorbox}[colback=defbox, colframe=defborder, arc=2pt,
  left=6pt, right=6pt, top=4pt, bottom=4pt,
  title={\small\textbf{Definition (Properties of Unary Operations)}},
  fonttitle=\small\bfseries]
\label{def:unary-laws}
Let $f : A \to A$ be a unary operation.
\begin{enumerate}[label=(\roman*)]
  \item $f$ is \emph{idempotent} if applying $f$ twice gives the same
        result as applying it once:
        \[
          \forall a \in A,\quad f(f(a)) = f(a).
        \]
  \item $f$ is an \emph{involution} if applying $f$ twice returns the
        original element:
        \[
          \forall a \in A,\quad f(f(a)) = a.
        \]
  \item $a \in A$ is a \emph{fixed point} of $f$ if
        \[
          f(a) = a.
        \]
\end{enumerate}
\end{tcolorbox}

\begin{remark}[Fully quantified form]
Idempotent: $(\forall a \in A)\; f(f(a)) = f(a)$.
\quad
Involution: $(\forall a \in A)\; f(f(a)) = a$.
\quad
Fixed point of $f$: a particular $a$ with $f(a) = a$.
\end{remark}

\begin{remark}[Involution is a strong form of bijectivity]
Every involution is a bijection: its own inverse.
$f \circ f = \mathrm{id}_A$ means $f^{-1} = f$.
Complex conjugation and negation on $\mathbb{Z}$ are both involutions.
An idempotent that is also a bijection must be the identity:
if $f \circ f = f$ and $f$ is invertible, compose both sides on the right
by $f^{-1}$ to get $f = \mathrm{id}_A$.
\end{remark}

% ---------------------------------------------------------
% Binary operation
% ---------------------------------------------------------

\begin{tcolorbox}[colback=defbox, colframe=defborder, arc=2pt,
  left=6pt, right=6pt, top=4pt, bottom=4pt,
  title={\small\textbf{Definition (Binary Operation)}},
  fonttitle=\small\bfseries]
\label{def:binary-op}
A \emph{binary operation} on $A$ is a function
\[
  \star : A \times A \to A.
\]
For $a, b \in A$ we write $a \star b$ for the image of $(a, b)$ under
$\star$.
\end{tcolorbox}

\begin{remark}[Fully quantified form]
$\star$ is a binary operation on $A$
$\;\iff\;$
$(\forall a, b \in A)\;(\exists!\, c \in A)\; a \star b = c$.
\end{remark}

\begin{remark}[Closure is definitional]
The codomain of $\star$ is $A$ itself, so $a \star b \in A$ for all
$a, b \in A$.
This is closure, and it is not an axiom — it is the meaning of
``$\star$ is a function $A \times A \to A$''.
\end{remark}

% ---------------------------------------------------------
% Laws of binary operations
% ---------------------------------------------------------

\subsection{Laws of Binary Operations}

\begin{tcolorbox}[colback=gray!6, colframe=gray!40, arc=2pt,
  left=6pt, right=6pt, top=4pt, bottom=4pt,
  title={\small\textbf{Laws of Binary Operations — Quick Reference}},
  fonttitle=\small\bfseries]
\small
\begin{tabular}{l l l}
\toprule
\textbf{Law} & \textbf{Condition} & \textbf{Detail} \\
\midrule
Associativity    & $(a \star b) \star c = a \star (b \star c)$         & \hyperref[def:associativity]{↓} \\
Commutativity    & $a \star b = b \star a$                             & \hyperref[def:commutativity]{↓} \\
Distributivity   & $a \star (b \diamond c) = (a \star b) \diamond (a \star c)$ & \hyperref[def:distributivity]{↓} \\
Cancellation (L) & $a \star b = a \star c \Rightarrow b = c$           & \hyperref[def:cancellation]{↓} \\
Cancellation (R) & $b \star a = c \star a \Rightarrow b = c$           & \hyperref[def:cancellation]{↓} \\
\bottomrule
\end{tabular}
\end{tcolorbox}

\vspace{1em}

\begin{tcolorbox}[colback=defbox, colframe=defborder, arc=2pt,
  left=6pt, right=6pt, top=4pt, bottom=4pt,
  title={\small\textbf{Definition (Associativity)}},
  fonttitle=\small\bfseries]
\label{def:associativity}
A binary operation $\star$ on $A$ is \emph{associative} if
\[
  \forall a, b, c \in A,\quad (a \star b) \star c = a \star (b \star c).
\]
\end{tcolorbox}

\begin{remark}[Fully quantified form]
$(\forall a, b, c \in A)\;(a \star b) \star c = a \star (b \star c)$.
\end{remark}

\begin{remark}[Algebraic meaning]
Associativity says that bracketing does not matter in any product of three
elements.
By induction this extends: any well-formed bracketing of a finite product
$a_1 \star a_2 \star \cdots \star a_n$ gives the same result.
This is the property that makes a finite sequence of elements have an
unambiguous product — and is therefore the minimum condition for defining
$a^n$ or $na$ meaningfully.
\end{remark}

\begin{tcolorbox}[colback=defbox, colframe=defborder, arc=2pt,
  left=6pt, right=6pt, top=4pt, bottom=4pt,
  title={\small\textbf{Definition (Commutativity)}},
  fonttitle=\small\bfseries]
\label{def:commutativity}
A binary operation $\star$ on $A$ is \emph{commutative} if
\[
  \forall a, b \in A,\quad a \star b = b \star a.
\]
\end{tcolorbox}

\begin{remark}[Fully quantified form]
$(\forall a, b \in A)\; a \star b = b \star a$.
\end{remark}

\begin{remark}[Independence from associativity]
Associativity and commutativity are logically independent.
Addition and multiplication on $\mathbb{Z}$ are both.
Matrix multiplication is associative but not commutative.
Subtraction on $\mathbb{Z}$ is neither.
Rock-paper-scissors (``beats'') is neither.
\end{remark}

\begin{tcolorbox}[colback=defbox, colframe=defborder, arc=2pt,
  left=6pt, right=6pt, top=4pt, bottom=4pt,
  title={\small\textbf{Definition (Distributivity)}},
  fonttitle=\small\bfseries]
\label{def:distributivity}
Let $\star$ and $\diamond$ be two binary operations on $A$.
\begin{enumerate}[label=(\roman*)]
  \item $\star$ \emph{distributes over} $\diamond$ \emph{on the left} if
        \[
          \forall a, b, c \in A,\quad
          a \star (b \,\diamond\, c)
          = (a \star b) \,\diamond\, (a \star c).
        \]
  \item $\star$ \emph{distributes over} $\diamond$ \emph{on the right} if
        \[
          \forall a, b, c \in A,\quad
          (b \,\diamond\, c) \star a
          = (b \star a) \,\diamond\, (c \star a).
        \]
  \item $\star$ \emph{distributes over} $\diamond$ if it distributes on
        both the left and the right.
\end{enumerate}
\end{tcolorbox}

\begin{remark}[Fully quantified form]
Left: $(\forall a,b,c \in A)\;
a \star (b \diamond c) = (a \star b) \diamond (a \star c)$.
Right: $(\forall a,b,c \in A)\;
(b \diamond c) \star a = (b \star a) \diamond (c \star a)$.
\end{remark}

\begin{remark}[Two operations required]
Distributivity is a relationship \emph{between} two operations on the same
set — it cannot be stated for a single operation alone.
When $\star$ is commutative, left and right distributivity are equivalent.
The paradigm example is multiplication distributing over addition on
$\mathbb{Z}$: $a(b + c) = ab + ac$.
This is the defining interaction between the two operations of a ring.
\end{remark}

\begin{tcolorbox}[colback=defbox, colframe=defborder, arc=2pt,
  left=6pt, right=6pt, top=4pt, bottom=4pt,
  title={\small\textbf{Definition (Cancellation Laws)}},
  fonttitle=\small\bfseries]
\label{def:cancellation}
A binary operation $\star$ on $A$ satisfies the
\begin{enumerate}[label=(\roman*)]
  \item \emph{left cancellation law} if:
        \[
          \forall a, b, c \in A,\quad
          a \star b = a \star c \;\Rightarrow\; b = c,
        \]
  \item \emph{right cancellation law} if:
        \[
          \forall a, b, c \in A,\quad
          b \star a = c \star a \;\Rightarrow\; b = c.
        \]
\end{enumerate}
$\star$ satisfies the \emph{cancellation laws} if it satisfies both.
\end{tcolorbox}

\begin{remark}[Fully quantified form]
Left: $(\forall a,b,c \in A)\;[a \star b = a \star c \Rightarrow b = c]$.
Right: $(\forall a,b,c \in A)\;[b \star a = c \star a \Rightarrow b = c]$.
\end{remark}

\begin{remark}[Cancellation and injectivity]
Left cancellation of $\star$ says: for each fixed $a \in A$, the map
$L_a : A \to A$ defined by $L_a(x) = a \star x$ is injective.
Right cancellation says the same for $R_a(x) = x \star a$.
In a group, cancellation follows from the existence of inverses.
In a finite cancellative structure, cancellation implies the existence of
inverses (Latin square property).
\end{remark}

% ---------------------------------------------------------
% Distinguished elements
% ---------------------------------------------------------

\subsection{Distinguished Elements}

\begin{tcolorbox}[colback=gray!6, colframe=gray!40, arc=2pt,
  left=6pt, right=6pt, top=4pt, bottom=4pt,
  title={\small\textbf{Distinguished Elements — Quick Reference}},
  fonttitle=\small\bfseries]
\small
\begin{tabular}{l l l}
\toprule
\textbf{Element} & \textbf{Condition} & \textbf{Detail} \\
\midrule
Left identity     & $e \star a = a$ for all $a$                     & \hyperref[def:identity-element]{↓} \\
Right identity    & $a \star e = a$ for all $a$                     & \hyperref[def:identity-element]{↓} \\
Identity element  & Both left and right identity                    & \hyperref[def:identity-element]{↓} \\
Left inverse of $a$   & $b \star a = e$                             & \hyperref[def:inverse-element]{↓} \\
Right inverse of $a$  & $a \star b = e$                             & \hyperref[def:inverse-element]{↓} \\
Inverse of $a$        & Both: $b \star a = a \star b = e$           & \hyperref[def:inverse-element]{↓} \\
Idempotent element    & $a \star a = a$                             & \hyperref[def:idempotent-element]{↓} \\
Absorbing element     & $a \star z = z \star a = z$ for all $a$     & \hyperref[def:absorbing-element]{↓} \\
\bottomrule
\end{tabular}
\end{tcolorbox}

\vspace{1em}

\begin{tcolorbox}[colback=defbox, colframe=defborder, arc=2pt,
  left=6pt, right=6pt, top=4pt, bottom=4pt,
  title={\small\textbf{Definition (Identity Element)}},
  fonttitle=\small\bfseries]
\label{def:identity-element}
Let $\star$ be a binary operation on $A$ and $e \in A$.
\begin{enumerate}[label=(\roman*)]
  \item $e$ is a \emph{left identity} for $\star$ if
        $\forall a \in A,\quad e \star a = a$.
  \item $e$ is a \emph{right identity} for $\star$ if
        $\forall a \in A,\quad a \star e = a$.
  \item $e$ is an \emph{identity element} for $\star$ if it is both a
        left identity and a right identity.
\end{enumerate}
\end{tcolorbox}

\begin{remark}[Fully quantified form]
Left identity: $(\forall a \in A)\; e \star a = a$.
Right identity: $(\forall a \in A)\; a \star e = a$.
Identity: both hold for the same $e$.
\end{remark}

\begin{remark}[Uniqueness of the identity]
If an identity element exists, it is unique.
Proof: suppose $e$ and $e'$ are both identities.
Then $e = e \star e' = e'$, where the first equality uses that $e'$ is a
right identity and the second uses that $e$ is a left identity.
\end{remark}

\begin{remark}[Left-only vs.\ two-sided]
A left identity and a right identity that are both present must be equal
(same argument as uniqueness).
A left identity alone need not equal a right identity alone —
but if both exist, they coincide, yielding the unique two-sided identity.
\end{remark}

\begin{tcolorbox}[colback=defbox, colframe=defborder, arc=2pt,
  left=6pt, right=6pt, top=4pt, bottom=4pt,
  title={\small\textbf{Definition (Inverse Element)}},
  fonttitle=\small\bfseries]
\label{def:inverse-element}
Let $\star$ be a binary operation on $A$ with identity element $e$,
and let $a \in A$.
\begin{enumerate}[label=(\roman*)]
  \item $b \in A$ is a \emph{left inverse} of $a$ if
        $b \star a = e$.
  \item $b \in A$ is a \emph{right inverse} of $a$ if
        $a \star b = e$.
  \item $b \in A$ is an \emph{inverse} of $a$ if
        $b \star a = a \star b = e$.
\end{enumerate}
\end{tcolorbox}

\begin{remark}[Fully quantified form]
Left inverse of $a$: $b \star a = e$.
Right inverse of $a$: $a \star b = e$.
Inverse of $a$: $b \star a = e$ and $a \star b = e$.
\end{remark}

\begin{remark}[Uniqueness in associative structures]
In an associative operation with an identity, if $a$ has both a left
inverse $b$ and a right inverse $c$, then $b = c$.
Proof: $b = b \star e = b \star (a \star c) = (b \star a) \star c = e \star c = c$.
Hence in a group, inverses are unique.
\end{remark}

\begin{tcolorbox}[colback=defbox, colframe=defborder, arc=2pt,
  left=6pt, right=6pt, top=4pt, bottom=4pt,
  title={\small\textbf{Definition (Idempotent Element)}},
  fonttitle=\small\bfseries]
\label{def:idempotent-element}
Let $\star$ be a binary operation on $A$.
An element $a \in A$ is \emph{idempotent} (with respect to $\star$) if
\[
  a \star a = a.
\]
\end{tcolorbox}

\begin{remark}[Fully quantified form]
$a$ is idempotent $\;\iff\; a \star a = a$.
\end{remark}

\begin{remark}[Examples]
In any structure with an identity $e$: $e \star e = e$, so the identity
is always idempotent.
In a Boolean algebra, every element is idempotent under both $\land$ and $\lor$.
In a ring, an idempotent $e$ satisfying $e^2 = e$ is called an
\emph{idempotent element} and plays a central role in decomposing the ring
into direct summands.
\end{remark}

\begin{tcolorbox}[colback=defbox, colframe=defborder, arc=2pt,
  left=6pt, right=6pt, top=4pt, bottom=4pt,
  title={\small\textbf{Definition (Absorbing Element)}},
  fonttitle=\small\bfseries]
\label{def:absorbing-element}
Let $\star$ be a binary operation on $A$.
An element $z \in A$ is \emph{absorbing} (or a \emph{zero element})
for $\star$ if
\[
  \forall a \in A,\quad a \star z = z \star a = z.
\]
\end{tcolorbox}

\begin{remark}[Fully quantified form]
$z$ is absorbing for $\star$
$\;\iff\;$
$(\forall a \in A)\;[a \star z = z \;\land\; z \star a = z]$.
\end{remark}

\begin{remark}[Uniqueness and the identity]
If an absorbing element exists, it is unique (same argument as for the
identity).
In a ring $(R, +, \cdot)$, the additive identity $0$ is absorbing for
multiplication: $a \cdot 0 = 0$ for all $a$.
An absorbing element and an identity element coincide only in the trivial
one-element structure: if $e$ is the identity and $z$ is absorbing,
then $e = e \star z = z$.
\end{remark}

% ---------------------------------------------------------
% Compatibility of relations and operations
% ---------------------------------------------------------

\begin{tcolorbox}[colback=defbox, colframe=defborder, arc=2pt,
  left=6pt, right=6pt, top=4pt, bottom=4pt,
  title={\small\textbf{Definition (Compatibility of a Relation with an Operation)}},
  fonttitle=\small\bfseries]
\label{def:compatibility}
Let $\star$ be a binary operation on $A$ and $R$ a binary relation on $A$.
We say $R$ is \emph{compatible with $\star$} if
\[
  a_1 \mathrel{R} a_2 \;\land\; b_1 \mathrel{R} b_2
  \;\Rightarrow\;
  (a_1 \star b_1) \mathrel{R} (a_2 \star b_2)
\]
for all $a_1, a_2, b_1, b_2 \in A$.
\end{tcolorbox}

\begin{remark}[Fully quantified form]
$(\forall a_1, a_2, b_1, b_2 \in A)\;
[a_1 \mathrel{R} a_2 \land b_1 \mathrel{R} b_2
\Rightarrow (a_1 \star b_1) \mathrel{R} (a_2 \star b_2)]$.
\end{remark}

\begin{remark}[What compatibility seeds]
Compatibility is the condition that connects the relational and operational
layers of algebraic structure.
When the relation is an equivalence relation, compatibility means the
operation descends to the quotient set — the equivalence relation is then
called a \emph{congruence relation}, and the quotient inherits the operation.
When the relation is an order, compatibility means the operation preserves
the order in both arguments — the condition defining ordered groups,
ordered fields, and lattices.
These structures are not named here; they are the subject of abstract algebra.
What is recorded is the single condition that makes them possible.
\end{remark}
